\section{Synchronization}


\begin{definition}
\textbf{Bad Interleaving:} Interleaving of operations from multiple threads that leads to incorrect program state.
\textbf{Critical Section:} Codeblocks that might lead to bad interleavings.
\end{definition}

\begin{definition}
\textbf{Thread Safety:} Nothing bad happening in all possible interleavings.
\end{definition}

\begin{definition}
\textbf{Thread Liveness:} Something good happening eventually.
\end{definition}

\subsection{Race Conditions}

\begin{definition}
\textbf{Race condition: } correctness depends on specific interleaving.
\end{definition}

\begin{definition}
\textbf{Low-level race condition: } insufficient synchronization of memory access (Data Races).
\textbf{High-level race condition: } wrong decomposition/algorithm (bad interleavings despite mutual exclusion).
\end{definition}

Generally we want one of three things when it comes to shared ressources:
\begin{itemize}
    \item \textbf{Thread local (Isolated Mutability):} We dont use the ressource of other threads.
    \item \textbf{Read only (Immutable):} We only read and do not change the ressource
    \item \textbf{Mutual exclusion (Synchronization):} We use the ressource of other threads but we ensure that only one thread accesses it at a time.
\end{itemize}

\subsection{Locks and Synchronized}

Every Java-Object has an \textbf{intrinsic lock} or lock monitor. These locks are \textbf{unique} and can be optaind by
threads which allows us to lock a certain codeblock. 

\begin{codeexample}
synchronized(lockObject) {
    i++;
}
\end{codeexample}

\begin{remark}
Synchronized can also be used as a keyword on a method to sync. the whole method.
\end{remark}

\begin{remark}
By default synchronized uses "this" as the lock-object and unlocks when an exception occurs.
\end{remark}

\begin{remark}
Java-Locks are reentrant meaning they can be locked multiple times by the same Thread recursively (It
will then have to release the look as many times as it aquired it).            
\end{remark}

Atomic integers have builtin synchronization, for example for incrementing/decrementing.

\begin{codeexample}
AtomicInteger counter = new AtomicInteger(0);
counter.incrementAndGet();
\end{codeexample}

\subsection{Wait, notify, notifyAll}

We can use wait() on a lock-object to tell the current thread to sleep until further notice if the lock is
not available. This allows use to save ressources.

\begin{codeexample}
synchronized(lockObject) {
    while(!condition) {
        lockObject.wait();
    }
    // do work
} 
\end{codeexample}

\begin{important}
Always wait inside a while loop, as this ensures checking the condition after beeing notified.
\end{important}

To wake one arbitrary or all waiting threads we can use notify() or notifyAll() on the lock-object.

\begin{codeexample}
synchronized(lockObject) {
    while (!condition) {
        lockObject.wait();
    }
    //do work
    notifyAll();
}
\end{codeexample}

\subsection{Lock API}

Instead of using built-in synchronized blocks we can also use locks from the java.util.concurrent.locks
package.

\begin{codeexample}
private Lock lock = new ReentrantLock();
...
void doWork(int x) {
    lock.lock();
    try {
        ...
    } finally {
        lock.unlock();
    }
}
\end{codeexample}

\begin{important}
Always use a \textbf{try-finally block to ensure the lock is released} even if an error occures.            
\end{important}

\begin{remark}
Use lock.tryLock(100, TimeUnit.MILLISECONDS) to accuire the lock only for
a given amount of time, avoiding deadlocks.
\end{remark}

{\scriptsize
\setlength{\tabcolsep}{2pt}
\begin{tabularx}{\columnwidth}{|l|X|X|}
\hline
Aspect & synchronized & Lock API \\\hline\hline
Handling & Auto (JVM) & Manual (lock/unlock) \\\hline
Release & Automatic & Manual (finally!) \\\hline
Scope & Block/Method & Flexible/Multi-method \\\hline
Waiting & Infinite block & tryLock() avoids block \\\hline
Interrupt & No & lockInterruptibly() ok \\\hline
Flex. & Low & High \\\hline
Comp. & Simple & Complex/Error-prone \\\hline
Risk & Lower & Higher (forget unlock) \\\hline
Usage & Simple sync & Advanced control \\\hline
\end{tabularx}
}

\subsection{Guidelines for Concurrent programming}

\begin{important}
\textbf{Thread-Locality:} Only share ressources if needed.
\end{important}

\begin{important}
\textbf{Immutability:} Use immutable data structures whenever possible.
\end{important}

\begin{important}
\textbf{Consitent locking:} Use locks to ensure mutual exclusion.
\end{important}

\begin{definition}
\textbf{Lock granularity } refers to the amount of code that is proteted by a lock in ways:
\begin{itemize}
        \item Size of synchronized blocks
        \item Number of objects a lock is used for
\end{itemize}
\end{definition}

\begin{theorem}
\textbf{Coarse-grained locking} is easier to implement but leads to less concurrency. \textbf{Fine-grained locking} is harder to implement without creating bugs but leads to more concurrency.
\end{theorem}

\begin{important}
Start with coarse-grained (simpler) and move to fine-
grained (performance) only if contention on the coarser locks becomes an issue.
\end{important}

\begin{definition}
\textbf{Atomicity: } An operation is atomic if it is performed as a single, indivisible unit. 
\end{definition}

\begin{important}
    Do not do expensive computations or I/O in critical
    sections, but also don't introduce race conditions
\end{important}
